A safe monitor state must express a declared requirement meaning. Residual soundness (RS) requires every accepted continuation to be safe after every admitted activation history; equality is exact, while strict inclusion can reject feasible updates. An empty history scope passes vacuously, and synthesis does not establish scope completeness or intent. Technical Appendix~\ref{ta:residual_formal} gives the full definition.

For Cell's start-only reference, initializing inspection as clear is exact, even while $B$ holds an uninspected workpiece. Initializing as pending is conservative for that reference. The duty actually needed at activation is therefore specified and proved separately:

\paragraph{An exact activation obligation with ordinary work.}\label{sec:cell-activation-duty}
For Cell, define an owed-inspection bit at activation to be 1 exactly when $B$ holds the workpiece. Thereafter $m_B$ sets it, $j$ clears it, and $m_A$ is forbidden while it is 1; other labels preserve it. Encoding 0 by $c$ and 1 by $d$ makes every non-error monitor transition identical to this obligation's rule, and precisely the forbidden departure enters error. Initialization and induction on the active suffix therefore establish exact enforcement of this activation obligation. In Figure~\ref{fig:cell-story}, the first path starts with bit 1 and ordinary inspection clears it; the second starts with bit 0. This is an explicit duty beginning at activation, stronger than the start-only reference scope above. It does not extend the 138 A checks to actual histories or to other requirements.

A Boolean event-history variable, or \emph{fluent}, records setting/resetting labels. The evaluation distinguishes three history scopes:
\par\noindent\textbf{A (activation-safe), for new requirements (NEW):} admit only safe pre-activation histories consistent with the observed component state.
\par\noindent\textbf{E (entry-scoped), for interval requirements (UPD):} start the reference at update entry.
\par\noindent\textbf{H (history-inclusive):} count pre-activation violations too; not claimed. Synthesis does not establish whether the chosen scope matches intent (S3).
